\section{La productividad es el producto: volante de datos y entrega real}

La sección anterior explicó que la empresa debe evitar dos atolladeros, el del relato y el de la rentabilidad; esta sección responde una pregunta más: si no se apoya en una historia, ¿qué vende en realidad la inteligencia corporizada? Los modelos grandes pueden decir «la inteligencia es el producto», porque el usuario compra directamente capacidades de conversación, redacción, código, búsqueda y Agent; la inteligencia corporizada, en cambio, tiene que decir «la productividad es el producto». Un robot no solo exhibe inteligencia: tiene que crear en el mundo físico una sustitución de trabajo, una mejora de eficiencia o una capacidad productiva nueva que puedan medirse. Lo que el cliente compra al final no es que «se parezca a una persona», sino la producción estable de tareas como acomodar anaqueles, cargar y trasladar, seleccionar artículos, inspeccionar, limpiar y repartir.

Esto también explica por qué la recolección de datos reales es tan cara. Contratar personas para teleoperar robots y recolectar datos reales tiene un costo alto, avanza despacio y cubre pocos escenarios; la descripción del episodio menciona que los datos reales representan solo el 1\% de los datos de entrenamiento y que el pipeline de datos sintéticos lleva el peso principal. Esa proporción no puede trasladarse mecánicamente a todas las empresas, pero muestra que, sin datos sintéticos, simulación y generalización dentro del escenario, el costo de los datos reales difícilmente sostendría la escala temprana de la inteligencia corporizada.

\begin{figure}[H]
\centering
\includegraphics[width=0.96\textwidth]{productivity-as-product.png}
\caption{La productividad es el producto: los modelos grandes venden inteligencia; la inteligencia corporizada tiene que vender productividad medible. Diagrama conceptual propio, elaborado a partir de la conversación de 01:55:17--02:13:51 y la descripción del video.}
\end{figure}

\begin{knowledgebox}{Lectura de la figura: de la tarea al ingreso hay que pasar por el costo}
En la figura, el paso de Task a Robot es solo el inicio; lo decisivo son Cost, Output y Revenue. Para que un robot genere valor de producto, tiene que bajar el costo por tarea, entregar una producción estable y verificable, y lograr que el cliente quiera volver a comprar. Una inteligencia que no reduce costos es solo una demostración cara.
\end{knowledgebox}

\subsection{Las cuentas de la teleoperación}

Esta subsección pone en números la afirmación «los datos reales son caros». La estimación que da Wang He (王鹤) es la siguiente: fabricar un robot humanoide full-size cuesta por lo menos del orden de cien mil yuanes; si se compran diez mil unidades para recolectar datos, la inversión en hardware es del orden de mil millones de yuanes. Más cara aún es la operación: si cada robot trabaja en dos turnos, incluso con dos personas teleoperándolo por turno, y se suman la anotación, el control de calidad y el mantenimiento, sostener durante un mes un sistema de recolección de diez mil unidades cuesta del orden de cientos de millones e incluso mil millones de yuanes. Estas cuentas explican por qué «depender por completo de datos reales de teleoperación» difícilmente puede ser la ruta general en la etapa temprana.

\begin{knowledgebox}{Las cuentas: por qué no basta con teleoperar para recolectar datos}
\begin{tabularx}{\textwidth}{p{0.22\textwidth}p{0.28\textwidth}X}
\toprule
Rubro de costo & Orden de magnitud aproximado & Lección \\
\midrule
Hardware de los robots & Diez mil unidades $\times$ del orden de cien mil yuanes & Antes de recolectar ya se requiere una inversión en activos del orden de mil millones de yuanes. \\
Personal de teleoperación & Varios turnos y varias personas por unidad & El personal no es un costo único, sino un burn continuo. \\
Anotación y control de calidad & Hay que revisar trayectorias, estados, fallas y resultados de la tarea & Los datos de acción no pueden etiquetarse a grandes trazos como una imagen común. \\
Escenario y mantenimiento & Instalaciones, reparaciones, reabastecimiento, manejo de excepciones & El mundo físico produce sin cesar problemas que no son del modelo. \\
\bottomrule
\end{tabularx}
\end{knowledgebox}

Esta es también una diferencia importante entre la conducción autónoma y la inteligencia corporizada. Un auto puede venderse al usuario, y el usuario genera datos de paso al manejarlo todos los días, de modo que el uso del producto absorbe el costo de los datos; si un robot todavía no puede crear valor por sí mismo, no hay manera de disfrazar la recolección de datos como comercialización. Vender al cliente un «robot sin funciones» para que él mismo recolecte y entrene puede generar ingresos a corto plazo, pero a largo plazo traslada al cliente el riesgo de investigación y desarrollo.

\subsection{La recipe de datos y el retorno desde las unidades entregadas}

Esta subsección convierte el problema de los datos en un problema de volante. La recipe de datos de la inteligencia corporizada puede entenderse así: los datos de internet y visuales aportan sentido común y representaciones; los datos sintéticos aportan la cola larga y un entrenamiento controlable; una cantidad pequeña de datos reales aporta grounding, y las fallas reales después de la entrega aportan un retorno continuo. En la etapa temprana, sin volumen de unidades entregadas, la proporción de datos reales puede ser muy baja; en cuanto se entra a escenarios reales, el retorno de datos conecta el producto con el pipeline de entrenamiento.

\begin{figure}[H]
\centering
\includegraphics[width=0.92\textwidth]{data-recipe-robotics.png}
\caption{Recipe de datos para sistemas corporizados: 1\% de datos reales, pipeline sintético, retorno desde las unidades entregadas y volante de datos. Diagrama conceptual propio, elaborado a partir de la conversación de 01:55:17--02:13:51 y la descripción del video.}
\end{figure}

\begin{knowledgebox}{Lectura de la figura: cuanto más arriba, más cerca de la realidad y más escaso}
La capa inferior de datos generales construye las capacidades básicas; la capa intermedia de datos sintéticos amplía la distribución de tareas; una cantidad pequeña de datos reales calibra, y el retorno desde las unidades entregadas aporta los casos de falla y de frontera más importantes. Al leer la figura hay que evitar un malentendido: que los datos reales sean una proporción baja no significa que sean poco importantes; precisamente porque son escasos, deben usarse para calibrar y validar todo el sistema.
\end{knowledgebox}

\begin{warningbox}{Un robot sin funciones no puede venderse como volante de datos}
La descripción de la entrevista menciona un fenómeno tricky: vender a otros un robot sin funciones, dejar que el cliente recolecte y entrene por su cuenta, y prometerle que en el futuro se convertirá en un empleado. Este modelo es muy riesgoso, porque traslada al cliente el riesgo de investigación y desarrollo y además agota la credibilidad de la industria.
\end{warningbox}

\begin{importantbox}{La definición de productividad}
Wang He define la productividad de forma muy sencilla: el trabajo realizado por unidad de tiempo debe equipararse al de una persona o, por lo menos, mejorar la eficiencia dentro de un límite claro. Si el robot es mucho más lento que una persona, no aguanta mucho tiempo, falla con frecuencia o necesita mucho respaldo humano, no es una productividad de calidad; incluso puede ser una productividad atrasada.
\end{importantbox}

\subsection{Resumen de la sección}

El producto de la inteligencia corporizada no es «parecer inteligente», sino productividad medible. El volante de datos también tiene que apoyarse en entregas reales: los datos sintéticos reducen el costo inicial, el retorno desde las unidades entregadas aporta calibración real y el valor para el cliente comprueba si la ruta funciona.
